\chapter{Software Architecture Design}
\label{chap:software-architecture-design}

\section{System Architecture Overview}
\label{section:architecture-overview}

The system is designed as a multi-modal AI assistant architecture that integrates mobile application interfaces, AI processing pipelines, and backend services to support communication for users with hearing and speech impairments.

\begin{figure}[h]
    \centering
    \includegraphics[width=0.8\textwidth]{system_architecture_diagram.png}
    \caption{System Architecture Diagram of TermTem System}
\end{figure}

The architecture follows a client-server model consisting of three main layers:

\begin{itemize}
    \item \textbf{Frontend (Mobile Application - Flutter)} \\
    Handles user interaction including text input, speech input, and real-time video capture for sign language recognition. It also renders outputs in text, speech, and sign language animation.
    
    \item \textbf{Backend Services (Go-based API Server)} \\
    Acts as the central orchestrator, managing communication between frontend and AI components, handling request routing, and coordinating data processing pipelines.
    
    \item \textbf{AI Processing Pipeline} \\
    A modular pipeline responsible for transforming user input into meaningful responses. It consists of:
    \begin{itemize}
        \item MediaPipe for extracting skeletal joint coordinates from video
        \item LSTM-based model for sign language recognition
        \item Large Language Model (LLM) for generating responses
        \item Output modules including Text-to-Speech and Sign Language Animation Generation
    \end{itemize}
\end{itemize}

The system supports multi-modal input (text, speech, video) and multi-modal output (text, speech, sign language video), making it accessible and flexible for users with different communication needs.

% --------------------------------------------------

\section{Software Design}
\label{section:software-design}

The system runtime behavior is centered around processing user input through a multi-stage pipeline and generating appropriate responses based on the selected output modality.

One of the most critical use cases is real-time sign language translation and response generation, which involves multiple interacting components.

The process begins when the user provides input via real-time video. The frontend captures the video stream and sends it to the backend. The backend forwards the video frames to MediaPipe, which extracts skeletal joint coordinates representing hand movements.

These coordinates are then processed by an LSTM-based model trained for Thai Sign Language recognition, producing a corresponding textual representation.

The recognized text is forwarded to the LLM, which generates a context-aware response. Depending on the selected output mode, the response follows one of three paths:

\begin{figure}[h]
    \centering
    \includegraphics[width=0.8\textwidth]{sequence_diagram.png}
    \caption{Sequence Diagram of TermTem System}
\end{figure}

\begin{enumerate}
    \item \textbf{Text Response} \\
    The generated text is sent directly back to the frontend.
    
    \item \textbf{Speech Response} \\
    The text is passed to the Google Text-to-Speech API to generate audio output.
    
    \item \textbf{Sign Language Video Response} \\
    The text is transformed into sign language grammar using an ASL Translator module. The processed sentence is split into tokens, which are mapped to corresponding sign language video clips stored in a database. These clips are sequenced and rendered into a continuous animation using Unity before being sent back to the frontend.
\end{enumerate}

This design ensures modularity and extensibility, allowing future improvements in individual components without affecting the entire system.

% --------------------------------------------------

\section{AI and Data Design}
\label{section:ai-data-design}

\subsection{Data Sources and Collection}

The system utilizes a combination of public, scraped, and self-collected datasets to train and improve the performance of the sign language recognition model.

\begin{enumerate}
    \item \textbf{Public Dataset (Kaggle)} \\
    The primary dataset consists of over 1,000 Thai sign language videos obtained from Kaggle. These videos are sourced from news agencies and have been pre-labeled and annotated with gesture information. Due to the availability of annotations, minimal preprocessing is required.
    
    \item \textbf{Web-Scraped Dataset} \\
    Additional data is collected from the Deaf Thai Association database. This dataset contains isolated sign language word videos. Since the data is not in a model-ready format, each video is processed using MediaPipe to extract joint coordinate sequences before being used for training.
    
    \item \textbf{Self-Collected Dataset (Future Work)} \\
    If the model performance is insufficient, additional data will be collected manually. This involves recording sign language videos and processing them through MediaPipe to generate training features.
\end{enumerate}

\textbf{Data Preprocessing includes:}
\begin{itemize}
    \item Video normalization (frame resizing and frame rate alignment)
    \item Keypoint extraction using MediaPipe
    \item Sequence padding and truncation for LSTM input
    \item Label encoding for classification
\end{itemize}

% --------------------------------------------------

\subsection{Model Design}

The core AI model for sign language recognition is based on a Long Short-Term Memory (LSTM) architecture.

LSTM is chosen due to its effectiveness in handling sequential temporal data, which is essential for interpreting dynamic hand gestures in sign language. It can capture long-term dependencies between frames, enabling more accurate gesture recognition.

The model pipeline is structured as follows:

\begin{itemize}
    \item \textbf{Input:} Sequence of joint coordinates extracted from MediaPipe
    \item \textbf{Model:} LSTM layers for temporal feature extraction
    \item \textbf{Output:} Predicted text label representing the recognized sign
\end{itemize}

\textbf{Model Architecture:}
\begin{itemize}
    \item Input shape: $(T, N)$ where $T$ is the number of frames and $N$ is the number of keypoint features
    \item LSTM Layer (128 units, return sequences=True)
    \item Dropout Layer (rate=0.2)
    \item LSTM Layer (64 units, return sequences=False)
    \item Dense Layer (64 units, ReLU activation)
    \item Dropout Layer (rate=0.2)
    \item Output Dense Layer (Softmax activation)
\end{itemize}

\begin{figure}[h]
    \centering
    \includegraphics[width=0.8\textwidth]{lstm_architecture_diagram.png}
    \caption{Architecture Diagram of Long Short-Term Model}
\end{figure}

For response generation, a Large Language Model (LLM) is integrated via API to generate context-aware conversational responses.

For sign language output generation, the system uses:
\begin{itemize}
    \item ASL Translator for grammar transformation
    \item Text Splitter for tokenization
    \item Token-to-video mapping using a database
    \item Unity engine for animation rendering
\end{itemize}

\textbf{Training Configuration:}
\begin{itemize}
    \item Loss Function: Categorical Crossentropy
    \item Optimizer: Adam
    \item Evaluation Metric: Accuracy
\end{itemize}

% --------------------------------------------------

\subsection{Evaluation}

The system performance is evaluated using both quantitative and qualitative metrics:

\begin{enumerate}
    \item \textbf{Recognition Accuracy} \\
    Measures the correctness of the sign language recognition model using metrics such as accuracy, precision, recall, and F1-score.
    
    \item \textbf{System Latency} \\
    Measures the response time from user input to system output, ensuring near real-time performance.
    
    \item \textbf{User Usability Evaluation} \\
    Conducted through surveys and interviews with users with communication impairments to evaluate usability, satisfaction, and overall user experience.
\end{enumerate}